\subsection*{What makes work avoidable}

This subsection is an existence argument depending on no measurement: a gap must exist between what
a dense-bitmap kernel \emph{spends} and what a cheaper but equally exact algorithm on the same
operands \emph{already achieves}, and representations whose cost is a function of structure rather
than of the universe can enter it. Both compared quantities are achieved; neither lower-bounds the
work the problem demands, and no such bound is claimed here. Every representation is exact and
interconvertible --- \rB, \rS, \rR, \rW\ and the complement view \rC\ encode the same set and every
cell returns the same cardinality --- so nothing is approximated and only the cost changes. Every
quantity in this subsection is a count of stored items or of operations, never a time; constants that convert
these counts to time are empirical and do not enter the derived results. Each result is stated here with its hypotheses; the proof, the worked
numerical checks and the secondary propositions that support but do not themselves state a result are
in Supplementary Note~\ref{note:proofs-avoidance}, for the results up to and including the size bound, and Supplementary Note~\ref{note:proofs-choice}, for the representation-choice optimisation problem that follows it.

Write $A,B \subseteq [0,m)$, $a=|A|$, $b=|B|$, densities $d_A=a/m$, $d_B=b/m$, word width $w$, vector
width $V$ words per operation. Complementation is an involution on density, $d\mapsto 1-d$, so define
the \emph{effective sparsity} $s_X=\min(d_X,1-d_X)\in[0,\tfrac12]$ and $s=\min(s_A,s_B)$.

\paragraph*{What cardinality alone determines.}
Given only $a$ and $b$ --- the cheapest metadata any selector holds --- the answer is confined to
\begin{equation}
\max(0,\,a+b-m) \;\le\; |A \cap B| \;\le\; \min(a,\,b),
\label{eq:interval}
\end{equation}
and every integer in the range is realised (Supplementary Note~\ref{note:proofs-avoidance}). Its width is the quantity the rest
of this subsection turns on:
\begin{equation}
\min(a,b)-\max(0,a+b-m) \;=\; \min\,(\,a,\;b,\;m-a,\;m-b\,) \;=\; m\,\min(s_A,s_B) \;=\; m\,s .
\label{eq:width}
\end{equation}
The four terms are the sizes of the four lists one could enumerate --- $A$, $B$, $A^{c}$, $B^{c}$ ---
so the uncertainty is the size of the smallest, and one near-empty or near-full side collapses it
alone. Above density one half the floor lifts off zero at $(d_A+d_B-1)m$: at $d_A=d_B=0.51$, two per
cent of the universe must overlap.

\paragraph*{The fixed cost, and therefore the headroom.}
A dense bitmap stores $\lceil m/w\rceil$ words whatever it contains, so $\cell{\rB}{\rB}$ performs
$\lceil m/(wV)\rceil$ vector AND-plus-popcount operations per pair, independently of $a$, $b$, of
structure, and of the answer. Against that, suppose each operand $X$ is available as a bitmap and as a
sorted array of whichever of $X$, $X^{c}$ is smaller --- $\Theta(m\,s_X)$ of space, the choice made
from $a$, $b$, $m$ in $O(1)$; no rank index is needed for this result. Then $|A\cap B|$ is answered in
$O(m\,s)$ membership probes (four cases, one per term of Eq.~\eqref{eq:width}; Supplementary Note~\ref{note:proofs-avoidance}).
Counting descriptor items touched, and taking $w \mid m$,
\begin{equation}
\frac{\text{items touched by } \cell{\rB}{\rB}}{\text{items touched by the cheapest list-based cell}}
\;=\; \frac{\lceil m/w\rceil}{m\,s} \;=\; \frac{1}{w\,s},
\label{eq:headroom}
\end{equation}
which exceeds one whenever $s<1/w$ and grows without bound as either operand leaves density one half
--- the headroom, derived rather than observed, and a comparison of two \emph{achieved} costs, not a
lower bound on any algorithm's work. Two of the four cases need a complement view, which is why \rC\
is part of the design rather than an afterthought for the dense tail. Charging measured per-item
costs --- a random probe against a vector word --- turns it into a crossover in $s$ of
\nm{measured crossover in effective sparsity}.

\paragraph*{From descriptor length to operation count.}
Everything else here counts \emph{stored items}, a fact about representation size; the paper's claim
is about work. One identity carries that step. Let $A$ be a bitmap with an $O(1)$ rank index,
$\mathrm{rank}_A(i)=|A\cap[0,i)|$, and $B$ a run container with maximal runs
$\{[u_i,v_i)\}_{i=1}^{r}$. The runs partition $B$, so
\begin{equation}
|A \cap B| \;=\; \sum_{i=1}^{r}\big(\,\mathrm{rank}_A(v_i)-\mathrm{rank}_A(u_i)\,\big),
\label{eq:rank}
\end{equation}
in exactly $2r$ lookups --- independent of $m$, of $|A|$, of $|B|$, and of the run lengths
(Supplementary Note~\ref{note:proofs-avoidance}). A run of any length is priced as a run of length one, because the kernel never
looks inside it; symmetrically $\cell{\rR}{\rR}$ by merge costs $O(r_A+r_B)$ with no index. This is
what makes a run count a cost parameter and not merely a storage parameter, and it is the only bridge
between the two here: it reaches \emph{operation count}, not time, since per-operation constants still
differ across cells.

\paragraph*{What randomness would predict.}
Let each position be set independently with probability $d$, and write $\kappa(\rho)$ for the number
of stored items representation $\rho$ must touch. The bitmap and array entries are exact; the run and
fill-compressed entries are the leading-order forms of exact Bernoulli expectations (Supplementary Note~\ref{note:proofs-avoidance}), written in
$s=\min(d,1-d)$,
\begin{equation}
\mathbb{E}\,\kappa(\rB) = \Big\lceil \tfrac{m}{w} \Big\rceil, \quad
\mathbb{E}\,\kappa(\rS) = m\,d, \quad
\mathbb{E}\,\kappa(\rR) = m\,d(1-d) = m\,s(1-s), \quad
\mathbb{E}\,\kappa(\rW) = \tfrac{m}{w}\big[\,1-(1-s)^{2w}-s^{2w}\,\big],
\label{eq:expected}
\end{equation}
where $\kappa(\rR)$ counts maximal runs and $\kappa(\rW)$ literal words plus fill groups. None of
these expectations is new. The \rW\ row is the word-aligned expected-size model of Wu, Otoo and
Shoshani \cite{wu2006wah}, whose bracket is identical to the one above under their payload width of
$w-1$ bits per word rather than $w$; the symmetry about $d=\tfrac12$ and the resulting
incompressibility at that density are theirs as well. The \rR\ row is the classical expected run
count for a uniformly random arrangement \cite{mood1940runs}, and the interval bound is
the Fr\'echet--Boole inequality. The density--clustering parameterisation is also due to
\cite{wu2006wah}; the design space it spans has since been mapped empirically for Roaring, WAH and
tree-encoded bitmaps by Lang et al. \cite{lang2020teb}, who report that an uncompressed bitmap reads
faster than any of the three compressed formats over $16 \le f \le 128$ and $0.01 \le d < 1$, and
that they expect ``a performance-optimized implementation to dominate an even larger space'' than
their unoptimized baseline did. That is the headroom this paper measures, stated from the compressed
side. They are restated here in one notation because the pairing matrix needs them as
inputs, not because they are results of this paper. With the
complement view making the best of \rS\ and $\rC\!\circ\!\rS$ cost $m\,s$, every entry is constant in
$s$ or strictly increasing on $[0,\tfrac12]$ and symmetric about $d=\tfrac12$: under this null model the
U-shaped cost curve is a theorem --- about the model and the bitmap's flatness, not about any method.
Randomness gives run containers nothing either, since
$\tfrac12 m s\le\mathbb{E}\,\kappa(\rR)\le m\,s+1$, so choosing \emph{among} \{\rS,
$\rC\!\circ\!\rS$, \rR, \rW\} buys at most a constant on structureless data. \rB\ is the exception and
the exception is the point: $\kappa(\rB)$ is $\lceil m/w\rceil$ at every density, so the spread across
the full matrix is Eq.~\eqref{eq:headroom} and unbounded. It is the envelope, not each cell, that
follows $\Theta(\min(m/w,\,m\,s))$.

\paragraph*{The deviation, which is the contribution.}
Real data is not random, and the representations part company with Eq.~\eqref{eq:expected} in a way
density cannot express. At fixed density $d\le\tfrac12$ with $k=dm$, $\kappa(\rB)$ and $\kappa(\rS)$
are constants, while $\kappa(\rR)$ ranges over every integer in $[1,k]$ and $\kappa(\rW)$ from $O(1)$
to $\Theta(\min(k,m/w))$ (proved by explicit families in Supplementary Note~\ref{note:proofs-avoidance}). Writing
$\sigma(\rho,X)$ for predicted descriptor length at $X$'s density over actual, $\sigma$ is identically
one for \rB\ and \rS, at least $\max(d,\,1-d+\tfrac{1}{m})\ge\tfrac12$ for \rR, and unbounded above:
real rows can be arbitrarily cheaper than random rows of the same density and at most about twice as
expensive.

One impossibility follows, exactly. A selector reading density --- or cardinality, or any function of
them --- evaluates one fixed function on every row of that density, so by the range above it
mis-predicts $\kappa(\rR)$ by an unbounded factor on some rows at every density; a fixed representation
cannot exploit $\sigma$ either, having committed before $\sigma$ is knowable. That rules out a family
of selectors without choosing among the survivors, and two further readings do not follow: whether a
mis-predicting model costs more than selecting nothing depends on which cell it picks and that cell's
actual cost, which is measured (Fig.~\ref{fig:selection-cost}b), not derived; and a run or occupancy
count is $O(1)$ metadata that sees $\sigma$ perfectly and could feed a model. What is proved is that
selection must be gated on a statistic that sees structure rather than on cardinality --- structural
summary or direct kernel measurement is resolved by the measured selector comparison in
Fig.~\ref{fig:selection-cost}b.

\paragraph*{The null model is a floor, and it is the same floor three times.}
Independent placement at fixed density is the \emph{extremal} arrangement for every mechanism used
here, always in the same direction (proofs in Supplementary Note~\ref{note:proofs-avoidance}). For a representation,
Eq.~\eqref{eq:expected} sits within a factor of two of the worst achievable descriptor. For a zone map
of bin width $W_z$, scattering maximises the occupied-bin fraction, so the summary's null cost is an
upper bound on its cost and a lower bound on its value. For the prefix filter of the narrowed-contract
mode, the entire mechanism depends on the operands through one dimensionless quantity, and any ordering
that concentrates rare elements into prefixes makes real prefixes collide less often than scattered
ones of the same length. In each case the null is what the mechanism is worth on data with no
structure to find, real data deviates from it only in the mechanism's favour, and that deviation is
exactly the quantity the mechanism converts into avoided work.

\paragraph*{The mid-density regime selects bitmap throughput.}
At $s=\tfrac12$ the interval has width $m/2$ and excludes nothing, the cheapest list-based cell spends
$m/2$ probes against $\lceil m/w\rceil$ words, and Eq.~\eqref{eq:expected} offers neither \rR\ nor \rW\
an advantage: Eq.~\eqref{eq:headroom} falls below one. In the generic case the mid-density band belongs
to $\cell{\rB}{\rB}$, and a selector choosing it there is correct rather than defeated. This bounds what
cardinality determines absent structure, not what any kernel must spend --- call it a \emph{metadata}
or \emph{identifiability floor}, never an information-theoretic one --- and a structured row at exactly
this density is a counterexample to the unqualified form (Supplementary Note~\ref{note:proofs-avoidance}).
Zone-map occupancy remains informative in this band because it records arrangement that cardinality
does not.

\paragraph*{Why the two mechanisms are selected between rather than stacked.}
Composing a zone map with a representation is not the product of the two, because filtering
\emph{concentrates} what it passes on: conditioned on a bin of width $W_z$ being occupied its expected
local density is $d/p$, $p=1-(1-d)^{W_z}$, not $d$ --- no summary hands a kernel anything sparser than
one bit per bin (Supplementary Note~\ref{note:proofs-avoidance}). Charging the summary pass $c/W_z$ of a bitmap scan and taking
the operands independent, the composed cost relative to a plain bitmap is
\begin{equation}
\frac{c}{W_z} \;+\; p_A\,p_B\;\kappa\!\Big(\frac{d}{p}\Big),
\label{eq:compose}
\end{equation}
$\kappa\le1$ being the envelope of Eq.~\eqref{eq:expected} normalised against the bitmap. The
composition never costs more than the summary alone, yet is floored at $c/W_z$, below which a sparse
representation is cheaper outright --- three regimes, the middle one exactly $x e^{-x} > c/w$ in
$x = W_z d$ bits per bin, opening where the summary undercuts enumeration and closing where bins stop
being empty. This is the analytic form of the gating result of Fig.~\ref{fig:selection-cost}c, and it
models \emph{work}: it bounds where a filter has work to remove and says nothing about what one costs
when it removes nothing.

\paragraph*{Each mechanism's liability is two-sided, for unrelated reasons at the two ends.}
A summary is indexed by the universe while every compressed representation is indexed by the data, so
their costs must cross. Charging the summary $c/W_z$ of a bitmap scan as in Eq.~\eqref{eq:compose},
and a one-word-per-element representation $w\,d$, the low crossing is
\begin{equation}
d_{\mathrm{low}} \;=\; \frac{c}{w\,W_z} \;=\; 3.1\times10^{-5} \qquad (c=1,\; w=64,\; W_z=512),
\label{eq:dlow}
\end{equation}
below which reading the summary costs more than enumerating the operands outright (Supplementary Note~\ref{note:proofs-avoidance}). At the other end $p_Ap_B\to1$: no bin is empty, the summary proves nothing, and only its own
$c/W_z$ remains. A run- or fill-based representation crosses at its own density, wherever its
descriptor length falls below $m/W_z$, but that a crossing exists at all is generic. Prefix filtering
has the same shape in $y=p^2/m$: for $y\ll1$ the surviving fraction is $y$ itself, so pruning improves
quadratically as prefixes shorten, while for $y>1$ prefixes collide by default and the index earns
nothing. This is why no mechanism here is applied unconditionally, and why each gate is a two-sided
test rather than a threshold.

\paragraph*{A narrowed contract, and what two integers then settle.}
Every statement above returns the cardinality itself. A caller may instead ask a narrower question ---
report only the pairs with $J(A,B)\ge t$, or only those with $|A\cap B|\ge\tau$, or restrict attention
to operands whose size falls in a band --- and the narrowing is itself exploitable, because
Eq.~\eqref{eq:interval} already says what sizes permit:
\begin{equation}
J(A,B) \;=\; \frac{|A\cap B|}{|A\cup B|} \;\le\; \frac{\min(a,b)}{\max(a,b)},
\label{eq:sizebound}
\end{equation}
so $J(A,B)\ge t$ requires $\min(a,b)\ge t\,\max(a,b)$, and $|A\cap B|\ge\tau$ requires
$\min(a,b)\ge\tau$ (proof in Supplementary Note~\ref{note:proofs-avoidance}). The hypotheses are that $a,b>0$ and $t\in(0,1]$;
nothing is assumed about arrangement, about $m$, or about how either operand is stored. Both conditions
are necessary and neither is sufficient, so the test discards pairs and never reports one --- the
survivors still go to an exact kernel, at the cost of one comparison between two integers already held
as metadata, no element of either set read. This is the Swamidass--Baldi pruning bound
\cite{swamidass2007bounds} together with the size (length) filter of
\cite{arasu2006exact}, in the two-sided form standardised by \cite{mann2016empirical} and applied
jointly with prefix filtering to Tanimoto over binary vectors by \cite{anastasiu2016tapnn}, restated here to place it against the mechanisms above
rather than as a result of this paper.
The operational narrowing and complete thresholded construction are given in Supplementary
Note~\ref{note:threshold}.

\paragraph*{What that removes, and why a wider corpus helps it.}
Modelling a corpus by drawing cardinalities independently with $\log a$ uniform on $[0,\ln K]$ --- the
density $\Pr(a)\propto 1/a$, exactly the $1/i$ cardinality spectrum this paper's benchmark protocol
mandates (``Corpora, provenance and loaders''), not an approximation to it ---
$\min(a,b)/\max(a,b)\ge t$ becomes $|\log a - \log b|\le\gamma$, $\gamma=\ln(1/t)$, and for two
independent uniforms on an interval of length $L=\ln K$,
\begin{equation}
\Pr\big[\text{pair can qualify}\big] \;=\; 1-\Big(1-\frac{\ln(1/t)}{\ln K}\Big)^{2},
\qquad \ln(1/t)\le\ln K,
\label{eq:survive}
\end{equation}
with the absolute form $\big(1-\ln\tau/\ln K\big)^{2}$ (proof and worked values in Supplementary Note~\ref{note:proofs-avoidance}, which agree with a $10^{5}$-sample simulation to three decimals). Expanding for small $\gamma$
gives $2\gamma\!\int\!\phi^{2}$ for any density $\phi$ of $\log a$, so the governing quantity is the
\emph{effective log-spread} $1/\!\int\!\phi^{2}$ of the size distribution, of which
Eq.~\eqref{eq:survive} is the log-uniform case. Two hypotheses carry real weight. Cardinalities are
integers, and $a=b$ passes at every $t$, so the continuous form is optimistic where sizes are small.
And the favourable scaling in $K$ --- Eq.~\eqref{eq:survive} decreases monotonically in $K$ --- is a
property of the log-uniform family and not of heavy tails generally: under the alternative reading in
which \emph{ranked} sizes scale as $1/i$, the surviving fraction is exactly $1-t$, independent of $K$,
since the effective log-spread saturates at two nats however far the sizes are spread. The favourable
scaling in $K$ is therefore claimed for the spectrum this paper benchmarks, not for every skewed
corpus.

\paragraph*{The size bound does not multiply with the within-pair mechanisms.}
Acting at a coarser granularity is not sufficient for savings to compose. The size test decides
whether a pair is formed at all, while representation choice and the zone map act inside a pair that
is formed, so the granularity argument of Eq.~\eqref{eq:compose} predicts a product --- and the
product is not achieved, because both mechanisms read the same statistic. The survivors of
Eq.~\eqref{eq:sizebound} are the pairs with $a\approx b$, and every within-pair mechanism is cheapest exactly
when $a$ and $b$ are \emph{dissimilar}, since its cost tracks $\min(a,b)$ against a fixed
$\lceil m/w\rceil$. Conditioning raises the mean cost of a survivor by
\begin{equation}
\frac{\mathbb{E}\big[\min(a,b)\;\big|\;\min/\max\ge t\big]}{\mathbb{E}\,\min(a,b)}
\;=\; \frac{L^{2}\big(K(1-t)-\gamma\big)}{\gamma\,(2L-\gamma)\,(K-L-1)}
\;\xrightarrow[\;t\to1\;]{}\; \frac{\ln K}{2},
\label{eq:shortfall}
\end{equation}
and the composed work exceeds the product of the two savings by that same factor; the derivation, the
numerical verification and the decorrelation control that isolates the shared statistic as the cause
are in Supplementary Note~\ref{note:proofs-avoidance}. The composition is still strongly worth taking. What is
refuted is only the product; the general rule is that mechanisms compose multiplicatively when they act
on different waste \emph{and} are driven by different operand statistics, and sharing a statistic costs
the composition Eq.~\eqref{eq:shortfall}'s factor even across granularities. Values here are analytic,
computed over the model above at $m=10^{6}$, and count work rather than time.

\paragraph*{Choosing a representation is a different problem from choosing a size, and only one of
them is separable.}
Everything above prices representations; a deployed library must \emph{choose} between them, and the
choice has to be stated as an optimisation before it can be called optimal. Partition the universe into
chunks of $M$ bits, each stored in one representation, $n=M/w$ words to a chunk and $e$ bits to a
stored array element. Given a corpus, a distribution $\pi$ over the chunk pairs an application
evaluates, and an operation $\omega$, the problem is to choose an assignment $\rho$ of representations
to chunks minimising $\sum_{(x,y)\sim\pi} C_\omega(\rho(x),\mu_x;\rho(y),\mu_y)$ subject to a byte
budget, where $\mu$ is the metadata of Eq.~\eqref{eq:expected} and $C_\omega$ counts items touched at
measured per-item constants. Storage is a \emph{separable} objective --- minimised one container at a
time --- whereas the query cost is quadratic in the assignment, since there is no such thing as the
cost of a container, only the cost of a pair (Supplementary Note~\ref{note:proofs-choice}). And the incumbent's three
container decisions --- array against bitset, array against run, bitset against run --- are one rule,
$\arg\min$ over serialized size, so its array-to-bitset threshold is $\tau_{\mathrm{stor}} = M/e$, a
property of the format carrying no machine constant and no property of the workload: the exact optimum
of a separable objective, and why it can be a compile-time constant. What follows is what changes
when the objective is not separable.

\paragraph*{The crossover is a formula, and the two constants in the field are its two limits.}
Write $c_w$ for a vectorised word \texttt{AND}-plus-popcount, $c_p$ for a membership probe into a
bitmap and $c_m$ for a merge step, and $\theta_{pw}=c_p/c_w$, $\theta_{mw}=c_m/c_w$. Promotion of a
chunk of cardinality $k$ from array to bitset is favourable against a bitmap partner exactly above
\begin{equation}
\tau_{\cap} \;=\; \frac{c_w}{c_p}\cdot\frac{M}{w}\;=\;\frac{n}{\theta_{pw}},
\qquad\text{whence}\qquad
\frac{\tau_{\mathrm{stor}}}{\tau_{\cap}} \;=\; \frac{w}{e}\,\theta_{pw},
\label{eq:tauratio}
\end{equation}
and against an array partner above $(\theta_{pw}/\theta_{mw}-1)\,k_Y$ for a partner of size $k_Y$
(proof in Supplementary Note~\ref{note:proofs-choice}). Pricing a byte at $\lambda\ge0$ and minimising query cost plus
$\lambda\,\mathrm{bytes}$ gives a threshold that runs from $\tau_{\cap}$ at $\lambda=0$ to
$\tau_{\mathrm{stor}}$ as $\lambda\to\infty$: every value between the two is optimal at some price of a
byte, and $\tau_{\mathrm{stor}}=M/e$ is the unique threshold at which promotion is never paid for in
bytes. The threshold also carries the pair count, so under any byte budget the optimum is not a
function of the container at all. The measured input is the symmetric crossover \nm{cardinality at
which an array--array merge stops beating a bitmap--bitmap popcount}, which fixes $\theta_{mw}$;
substituted into Eq.~\eqref{eq:tauratio} at $M=2^{16}$, $w=64$, $e=16$ together with \nm{measured ratio
of a bitmap probe to a merge step} it gives \nm{the ratio of the two thresholds}, recovering the
divergence between the two objectives from measurements of the constants and no free parameters. Both
crossovers share $c_w$, so measuring them together over-determines the model and the probe-to-merge
ratio becomes a checkable prediction rather than an input.

\paragraph*{A fixed threshold is optimal, under five hypotheses, and each of them fails.}
The loose form is false, so state the hypotheses first: if the candidates are
only \rS\ and \rB, the operation is a count-only intersection, the per-item constants do not depend on
the resident footprint, there is no byte budget, and every chunk faces the same partner mix, then the
optimal assignment \emph{is} the fixed threshold $\tau_{\cap}$ of Eq.~\eqref{eq:tauratio} --- a machine
constant, independent of the corpus. Dropping any one hypothesis breaks it, with an exact witness for
each (Supplementary Note~\ref{note:proofs-choice}): admitting \rR\ reintroduces the impossibility above, specialised to the
container decision, since $\cell{\rR}{\rB}$ costs a strictly increasing function of run count while
$\cell{\rS}{\rB}$ does not depend on it; letting the partner mix vary makes the optimal threshold
\begin{equation}
\tau^{*}(\beta,\bar k)\;=\;\frac{\beta\,c_w n+(1-\beta)(c_p-c_m)\,\bar k}
                                 {\beta\,c_p+(1-\beta)\,c_m},
\label{eq:taustar}
\end{equation}
where $\beta$ is the pair-weighted fraction of a chunk's partners that are bitmaps and $\bar k$ the
mean cardinality of the array partners: a statistic of the workload rather than of the machine, and
self-referential --- an optimal fixed
threshold is a fixed point $\tau=\tau^{*}(\beta(\tau),\bar k(\tau))$, which exists but need not be
unique; admitting a byte budget makes the threshold its dual variable, sweeping the whole interval
above; changing the operation flips the matrix's sign, since cardinality queries have no operation
dimension at all (union, difference and symmetric difference are inclusion--exclusion over the
intersection) but a \emph{materialising} union with a bitmap operand costs $n$ words whatever the other
side is, so no single stored assignment serves a mixed workload. The fifth hypothesis --- that the
per-item constants do not depend on footprint --- is the one that is not analytic and the one the
measurements say dominates: $c_w$ is not a constant, footprint re-enters through it, and a work model
that prices the choice cannot price its consequence, the same divergence recorded under
Eq.~\eqref{eq:compose} and treated as a measurement question, not an analytic one.

\paragraph*{Metadata settles the decision wherever the decision matters.}
Two identifiability questions must not be merged. Eq.~\eqref{eq:width} says what cardinality
determines about the \emph{answer}, pinned only at the density extremes. What cardinality determines
about the \emph{decision} is different and more favourable: every cost in this paragraph and the two
before it is a closed form in $(k,r,\ell+g)$ and the constants, so a selector holding those three
integers per operand evaluates the ordering exactly, in $O(1)$, touching no operand data --- which is
what ``selection must be near-free'' means, and a consequence of the cost model rather than an
engineering aspiration (proof in Supplementary Note~\ref{note:proofs-choice}). Holding only $k$, the ordering is not
determined once \rR\ is a candidate; one extra integer restores it and more cardinality precision does
not. Suppose the constants are known within a factor $\eta\ge1$ --- the uncertainty of the modelled
costs, distinct from the $\gamma$ of the size bound above --- and let $\Lambda$ be the ratio of the two
cheapest modelled costs: if $\Lambda>\eta$ the ordering is determined; if $\Lambda\le\eta$ it is
not, and choosing the second-best costs at most $\eta$. So for every pair, either the metadata
settles the decision or the decision does not matter to within the precision of the constants, and the
undetermined set is an $\eta$-neighbourhood of the boundaries rather than a region of unbounded error.
The caveat that makes $\eta$ interesting: it is the uncertainty of the \emph{model}, and where the
model omits the memory system it is not small --- the derived reason to measure a candidate rather than
only to model it.

\paragraph*{Why an advantage survives repairing the kernel beneath it.}
A pairwise batch costs $C_{\mathrm{meta}}$ plus the cell cost summed over the pairs a summary does not
discard and the regions it does not exclude; those sets are functions of cardinality, bin occupancy and
the data, and of nothing else, so no change of representation moves either set (Supplementary Note~\ref{note:proofs-choice}).
Consequently, if a better kernel divides every cell cost by $\alpha$, the advantage over the same batch
without avoidance is invariant in $\alpha$ when the summary is free and erodes only through the
summary's share of the surviving cost: improving a kernel cannot subsume avoidance, because avoidance
changes which work is asked for and a kernel changes only what asked-for work costs. Which savings
multiply is settled by the same rule as Eq.~\eqref{eq:shortfall}: a zone map reads bin occupancy while
representation choice reads cardinality and run count, so those multiply; the size bound reads
cardinality, which representation choice also reads, so those do not, and a container-threshold gain
must not be multiplied into a size-bound gain. Quantitatively, against an incumbent already holding the
compute-optimal threshold of Eq.~\eqref{eq:tauratio}, the modelled advantage of adding a zone map of bin
width $W_z$ peaks at
\begin{equation}
G_{\max}\;=\;\frac{\theta_{pw}\,w}{2\sqrt{c\,W_z}}
\qquad\text{at}\qquad d^{*}=\frac{\sqrt{c}}{W_z^{3/2}},
\label{eq:gmax}
\end{equation}
valid on $(\theta_{pw}w\sqrt{c})^{2/3}\le W_z\le\theta_{pw}w$ and saturating at
$\sqrt{\theta_{pw}w/c}\,/2$ above it, and falling below one outside a window whose lower edge is
Eq.~\eqref{eq:dlow} with the per-item constants carried, $c/(\theta_{pw}wW_z)$ (proof in Supplementary Note~\ref{note:proofs-choice}). Eq.~\eqref{eq:gmax} is a \emph{null-model} ceiling on structureless data, and by the
one-sidedness of $\sigma$ above, real corpora can only fall below the null, so a measured advantage
exceeding it is evidence of structure and must be reported as such rather than as agreement. And
Eq.~\eqref{eq:gmax} excludes the narrowed contract entirely: Eq.~\eqref{eq:survive} supplies its own
factor, and by the rule above the two are not multiplied.

\paragraph*{The irreducible regime selects throughput.}
The mid-band above is not only where these results stop applying; it is where a different optimization
applies, and the two are complementary rather than competing. Where Eq.~\eqref{eq:headroom} falls below
one and $p_Ap_B\to1$, the work $\cell{\rB}{\rB}$ performs is, in the generic case, work that is going to
be performed, and what remains to optimize is the per-cycle throughput of its $\lceil m/(wV)\rceil$
vector \texttt{AND}-plus-popcount operations --- a quantity none of the results above speaks to, and one
the popcount literature has already driven close to the hardware limit \cite{mula2018popcount}.
Eq.~\eqref{eq:headroom} and Eq.~\eqref{eq:compose} locate the boundary between the two domains; they do
not rank them.
